\documentclass[11pt]{article}
\usepackage[a4paper,margin=1in]{geometry}
\usepackage{amsmath,amssymb}
\usepackage{graphicx}
\usepackage{booktabs}
\usepackage[hidelinks]{hyperref}
\usepackage{authblk}

\newcommand{\Fl}{\mathbb{F}_\ell}
\title{\textbf{Addendum: A Galois-Fluid Stress Test of the $\chi(3)$ Channel\\
\large Frobenius traces $a_\ell(E_p)$ confirm CRT orthogonality}}
\author{R.\ Chen}
\affil{GUT Geoservice Inc., Montr\'eal}
\date{\textit{6N Arithmetic Geodynamics series} $\cdot$ Addendum to the $\chi(3)$ Channel note $\cdot$ 2026}

\begin{document}
\maketitle

\noindent
The $\chi(3)$-channel principle predicts that the 6N wing ($p\bmod 3$) can reach an
$L$-function-derived invariant of $E_p$ only through the $q=3$ Euler factor. This addendum
subjects that prediction to a deliberately adversarial test, using a quantity drawn from a
\emph{different} Euler factor: the Frobenius trace $a_\ell(E_p)$ for $\ell = 7, 13$. If the wing
could influence the curve's local data at $\ell \neq 3$, the principle would have a leak. It does not.

\paragraph{Design.} For every prime $3 < p < 2\times10^4$ (2260 primes) we form $E_p : y^2 = x^3 - p$
and extract $a_7(E_p)$ and $a_{13}(E_p)$ via \texttt{E.ap($\ell$)} in SageMath. We record, alongside
the 6N wing, the residues $p\bmod 7$ and $p\bmod 13$. These residues are not noise to be stripped:
they are the \textbf{positive control}. The reduction $E_p \bmod \ell$ depends only on $-p \bmod \ell$,
so $a_\ell$ is a deterministic function of $p \bmod \ell$; the control \emph{must} fire if the
instrument works. The main test then asks whether the wing adds anything.

\paragraph{Result.} The two questions separate as cleanly as an experiment in this series ever has.

\begin{table}[h]
\centering
\begin{tabular}{lcc}
\toprule
 & controller & dependence of $a_\ell$ \\
\midrule
Positive control & $p \bmod \ell$ \ ($\ell = 7, 13$) & $\eta^2 = 1.000$ \ (fully determined) \\
Main test        & $p \bmod 3$ \ (the 6N wing)        & $|d| \approx 0.005$ \ (null) \\
\bottomrule
\end{tabular}
\end{table}

\noindent
The positive control is not approximate: each residue class $p \bmod 7$ maps to a single fixed value
of $a_7$ (and likewise $a_{13}$ across its twelve classes), so $p \bmod \ell$ explains $100\%$ of the
variance. The main test is uniformly null --- for $a_7$, $d = -0.006$; for $a_{13}$, $d = -0.004$; the
same for $|a_\ell|$, $a_\ell^2$, and the Sato--Tate angle $\theta = \arccos\!\big(a_\ell/2\sqrt{\ell}\big)$
(all $|d| < 0.008$).

\begin{figure}[h]
\centering
\includegraphics[width=0.92\textwidth]{fig_galois_fluid.pdf}
\caption{Left: the positive control fires perfectly --- $a_7$ is a fixed function of $p \bmod 7$
($\eta^2 = 1.000$). Right: the two wings' $a_7$ distributions coincide (Cohen $d = -0.006$). The same
quantity is fully determined by $p \bmod 7$ and entirely independent of $p \bmod 3$.}
\end{figure}

\paragraph{Interpretation: $a_\ell$ is a CRT coordinate.} The result is the cleanest possible display of
orthogonality. One and the same quantity $a_\ell$ is \emph{completely} determined by $p \bmod \ell$ yet
\emph{completely} independent of $p \bmod 3$. This is not the wing being blocked by a barrier; it is the
Chinese Remainder Theorem itself: $p \bmod 7$ and $p \bmod 3$ are independent coordinates, so a function
of one carries no information about the other. The probe was framed as ``decrepitating the CRT lattice to
extract a deep fluid,'' but the fluid it extracts --- $a_\ell$ --- \emph{is} a CRT coordinate: bombarding
$p$ at temperature $\ell$ yields exactly $p \bmod \ell$, an axis orthogonal to the wing.

\paragraph{Consequence for the principle.} This independently verifies the $\chi(3)$-channel principle on
its hardest case. The wing acts only in the $q=3$ factor; $a_7$ and $a_{13}$ live in the $q=7$ and $q=13$
factors; different doors, no leak. It also closes the family: any probe ``wing vs $a_\ell$'' for any
$\ell$ will return $\eta^2 = 1$ against $p \bmod \ell$ and $d = 0$ against the wing, because CRT makes all
the residue axes mutually orthogonal. No further $\ell$ need be drilled.

\paragraph{Archival note.} Data (\texttt{Galois\_Fluid\_Isotopes\_Log.csv}, 2260 primes), the extraction
script (\texttt{Galois\_fluid\_log.sage}), and the analysis (\texttt{Galois\_fluid\_analysis.py}) are
archived with this addendum at \url{https://github.com/Ruqing1963/6N-chi3-channel}. The positive control
($\eta^2 = 1.000$) certifies the measurement; the result carries the highest confidence of any negative
in the series, because the instrument is provably working and the null is exact.
\end{document}
